%===============================================================================
% Konfigurationsdatei für LaTeX-Template
%===============================================================================

%===============================================================================
% Geometry und Seitenformat
%===============================================================================
\usepackage[a4paper, left=3cm, right=3cm, top=2cm]{geometry}

%===============================================================================
% Sprache und Encoding
%===============================================================================
\usepackage[ngerman]{babel}
\usepackage[utf8]{inputenc}
\usepackage[T1]{fontenc}

%===============================================================================
% Grafiken und Bilder
%===============================================================================
\usepackage{graphicx}
\usepackage{pdfpages}
\usepackage{subcaption}
\usepackage{float}

%===============================================================================
% Grafiken mit TikZ
%===============================================================================
\usepackage{tikz}
\usepackage{pgfplots}
\pgfplotsset{compat=1.18}
\usetikzlibrary{calc, angles, quotes}

%===============================================================================
% Mathematik
%===============================================================================
\usepackage{amsmath}
\usepackage{amssymb}
\usepackage{amsfonts}

%===============================================================================
% Einheiten und Zahlenformatierung
%===============================================================================
\usepackage[range-units=repeat, multi-part-units=brackets, range-phrase={~bis~},
    separate-uncertainty, locale=DE, per-mode=symbol-or-fraction]{siunitx}

%===============================================================================
% Tabellen
%===============================================================================
\usepackage{booktabs}
\usepackage{array}
\usepackage{longtable}
\usepackage{multirow}

%===============================================================================
% Farben
%===============================================================================
\usepackage{xcolor}

%===============================================================================
% Code-Highlighting mit Minted
%===============================================================================
\usepackage[newfloat]{minted}
\usemintedstyle{tango}  % oder andere Stile: monokai, perldoc, default
\definecolor{bg}{rgb}{0.95,0.95,0.95}
\setminted{
    bgcolor=bg,
    fontsize=\small,
    linenos=true,
    breaklines=true,
    frame=lines
}

% Fallback für listings (wenn Minted nicht verfügbar)
\usepackage{listings}
\lstset{
    basicstyle=\ttfamily\small,
    breaklines=true,
    commentstyle=\color{gray},
    keywordstyle=\color{blue},
    stringstyle=\color{red},
    showstringspaces=false,
    frame=single,
    captionpos=b,
    numbers=left,
    numberstyle=\tiny\color{gray},
    backgroundcolor=\color{white!95!black}
}

%===============================================================================
% Zitate und Anführungszeichen
%===============================================================================
\usepackage{csquotes}

%===============================================================================
% Hyperlinks und PDF-Navigation
%===============================================================================
\usepackage[hidelinks]{hyperref}
\usepackage{bookmark}

%===============================================================================
% Kopf- und Fußzeile
%===============================================================================
\usepackage{fancyhdr}
\pagestyle{fancy}
\fancyhf{}
\rhead{\thepage}
\lhead{\workTitle}
\renewcommand{\headrulewidth}{0.4pt}

%===============================================================================
% Kapitel und Überschriften
%===============================================================================
\usepackage{titlesec}
\titleformat{\chapter}{\Huge\bfseries}{Kapitel \thechapter:}{20pt}{\Huge\bfseries}
\titleformat{\section}{\large\bfseries}{\thesection}{1em}{}
\titleformat{\subsection}{\normalsize\bfseries}{\thesubsection}{1em}{}
\titleformat{\subsubsection}{\normalsize\itshape}{\thesubsubsection}{1em}{}

%===============================================================================
% Abbildungen und Tabellen
%===============================================================================
\usepackage{caption}
\captionsetup[figure]{labelfont=bf, textfont=normalfont}
\captionsetup[table]{labelfont=bf, textfont=normalfont}

%===============================================================================
% Benutzerdefinierte Variablen
%===============================================================================
\def\workTitle{Titel der Arbeit}
\def\workAuthor{Dein Name}
\def\workDate{\today}
\def\workUniversity{Universität Name}
\def\workInstitute{Institut Name}

%===============================================================================
% Theoreme und Definitionen
%===============================================================================
\usepackage{amsthm}

\newtheorem{theorem}{Satz}[section]
\newtheorem{lemma}[theorem]{Lemma}
\newtheorem{corollary}[theorem]{Korollar}
\newtheorem{definition}[theorem]{Definition}
\newtheorem{remark}[theorem]{Bemerkung}

%===============================================================================
% Eigene Befehle und Makros
%===============================================================================

% Betonung für wichtige Texte
\newcommand{\highlight}[1]{\colorbox{yellow!20}{#1}}

% Inline-Code
\newcommand{\code}[1]{\texttt{#1}}

% Referenzen
\newcommand{\figref}[1]{Abbildung~\ref{#1}}
\newcommand{\tabref}[1]{Tabelle~\ref{#1}}
\newcommand{\secref}[1]{Abschnitt~\ref{#1}}
\newcommand{\chref}[1]{Kapitel~\ref{#1}}

%===============================================================================
% Mikrotypographie (Verfeinerung)
%===============================================================================
\usepackage{microtype}

%===============================================================================
% Listen und Aufzählungen
%===============================================================================
\usepackage{enumitem}
\setlist[itemize]{leftmargin=*}
\setlist[enumerate]{leftmargin=*}

%===============================================================================
% Todo-Notizen (optional)
%===============================================================================
\usepackage[colorinlistoftodos]{todonotes}
